\begin{figure}[tbp]
\centering
\begin{tikzpicture}[>=Stealth,font=\small]
 \begin{scope}
  \node[font=\small\bfseries] at (1.65,2.5) {Source heights};
  \fill[blue!5] (-.1,-.15) rectangle (2.75,2.12);
  \draw[gray,dashed] (2.75,-.2)--(2.75,2.15);
  \node[anchor=east] at (2.65,1.78) {$\alpha_i=s_i(\mathcal C)$};
  \draw[blue!65!black,very thick] (.55,.95)--(2.75,.95);
  \fill (.55,.95) circle (1.7pt);
  \fill (2.75,.95) circle (1.7pt);
  \node[below] at (.55,.95) {$p_0-\varepsilon v$};
  \node[right] at (2.75,.95) {$p_0$};
  \node[above] at (1.55,.95) {$E=p_0+E_0$};
  \node at (1.6,-.55) {$P(E)=P(E_0)>0$};
 \end{scope}
 \begin{scope}[xshift=7cm]
  \node[font=\small\bfseries] at (0,2.5) {Target heights relative to the pivot};
  \draw[blue!65!black,very thick,->] (0,.95)--(-2.05,.95);
  \draw[orange!80!black,very thick,->] (0,.95)--(2.05,.95);
  \fill (0,.95) circle (1.7pt);
  \node[below] at (0,.95) {$0$};
  \node[above] at (-1.1,.95) {$E'_1=A_1E_0$};
  \node[above] at (1.1,.95) {$E'_2=-A_1E_0$};
  \node at (0,-.55) {$P(E'_1\cup E'_2)=2P(E_0)$};
 \end{scope}
\end{tikzpicture}
\caption{The height geometry of the reflection contradiction. The two affine profiles agree at the facet pivot and send the inward segment to opposite target segments. Alternating bundles detect the volume exponent of their union. The picture is schematic: the exponential counts come from bracket-generated boxes, and the volume argument accounts for the nonuniqueness of lifts to $N\rtimes D$.}
\label{fig:reflection}
\end{figure}
